\section{Background on DP and MDPs\label{sec:mdp}}

DDC and IRL apply  the recursive mathematical characterization of optimal sequential decision
making under uncertainty known as  {\it dynamic programming\/} (DP), \citet{Bellman:1957}.\footnote{\footnotesize See
\citet{bertsekas:2017} for an excellent introduction to the DP literature.} We focus 
on stationary, infinite horizon Markovian Decision Problems (MDP), a subclass of DPs introduced by Bellman and \citet{Howard:1960} that 
are the dominant paradigm in SE and RL.
MDPs are defined by  1) a state space $S$, 2) an action space $A$, and  3) the objects  $\{r,p,\beta\}$, where $r(s,a)$ is the {\it reward or utility function\/}
providing the payoff to a decision maker (DM), $p(s'|s,a)$ is a {\it transition density\/} for next period's state when the state is $s$ and action is $a$
and $\beta \in (0,1)$ is a discount factor.  To simplify exposition and avoid technicalities, we assume
that  $S$ and $A$ are finite sets with $|S|$ and $|A|$ elements, respectively.
A policy or {\it decision rule\/} $\pi$ is a conditional probability of choosing 
a feasible action $a \in A(s)$ for all $s \in S$. It is called a {\it mixed strategy\/} in game theory  to 
distinguish it from a {\it pure strategy,\/} i.e. a deterministic function $a=\delta(s)$ that
chooses a single optimal action with probability 1. Since mixed strategies include pure strategies as a special case, 
we can express the DM's objective as the choice of a policy
$\pi$ that  maximizes expected discounted rewards given by
\begin{equation}
       V(s)= \max_\pi  V_\pi(s),  \quad \mbox{where} \;  V_\pi(s) = E_\pi \left\{\sum_{t=0}^\infty \beta^t r(s_t,a_t)| s_0=s\right\},
\label{eq:vdef}
\end{equation}
where $E_\pi$ denotes the expectation operator over the Markov process $\{s_t,a_t\}$ implied by the transition probability $p$ and policy $\pi$.
The optimal value function can be defined recursively via {\it Bellman's equation\/}
\begin{equation}
            V(s) = \max_{a \in A(s)}\left[ r(s,a)+ \beta \sum_{s' \in S} V(s')p(s'|s,a)\right] \equiv \Gamma(V)(s).
\label{eq:bellman}
\end{equation}
$V$ is the fixed point $V=\Gamma(V)$ to an operator $\Gamma(V)$ known as the ``Bellman operator'' 
which can be shown to be a {\it contraction mapping\/} 
(i.e. for any vectors $V,W \in R^{|S|}$ we have $\|\Gamma(V)-\Gamma(W)\| \le \beta \|V-W\|$ where $\beta$ is less than 1
and $\|\cdot\|$ is the Euclidean norm). This implies that a  unique solution $V$ to Bellman's equation exists, 
and the optimal policy $\pi$ takes the form of a pure strategy, i.e. a deterministic function $a=\delta^*(s)$
that can be recovered from $V$ by
\begin{equation}
       \delta^*(s) = \argmax_{a \in A(s)}\left[ r(s,a)+\beta \sum_{s' \in S} V(s')p(s'|s,a)\right],
\label{eq:bellman_policy}
\end{equation}
since generically a unique action $a^*=\delta^*(s)$ attains the max on the right hand side of (\ref{eq:bellman_policy}).
Thus, a mixed strategy $\pi$ cannot be optimal in general. This appears to challenge the relevance of the MDP
framework since most strategies we estimate/learn empirically will be mixed rather than pure. 
The DDC literature resolves this by assuming that  
we only observe a subset of the full state vector the individual observes. This allows an optimal
pure strategy to appear to us as a mixed strategy given our inability to observe the agent's full information set. 
It also facilitates statistical inference of the MDP as we show in Section 3.
We now summarize the key methods for solving for  $V$ and $\delta^*$.


\subsection{Successive Approximations}

Since a contraction mapping $\Gamma$ has a unique fixed point $V=\Gamma(V)$, the contraction property guarantees that 
the sequence $\{V_k\}$ given by $V_k=\Gamma(V_{k-1})$, for $k=1,2,\ldots$ starting from
any initial starting guess $V_0$ converges to the fixed point $V=\Gamma(V)$ at a geometric rate in $\beta$,
which is very slow when $\beta$ is close to 1. 

\subsection{Policy Iteration}
\label{section:pi}

The \citet{Howard:1960} {\it policy iteration\/} (PI) algorithm is a faster method for solving MDPs that 
alternates between {\it policy valuation\/} and {\it policy improvement\/} steps. At iteration $k$ the PI algorithm has a candidate
policy $\delta_k$. The policy valuation step calculates the value of $\delta_k$,  $V_{\delta_k}$, by solving the linear system
\begin{equation}
     V_{\delta_k} = r_{\delta_k} + \beta E_{\delta_k} V_{\delta_k},
\label{eq:policy_valuation}
\end{equation}
where $r_{\delta_k}$ is the vector whose $s^{th}$ element is $r(s,\delta_k(s))$ and $E_{\delta_k}$ is the $|S| \times |S|$ matrix whose
element in row $s$ and column $s'$ is given by $E_{\delta_k}(s,s')=p(s'|s,\delta_k(s))$. For any feasible deterministic policy $\delta$
the matrix $E_{\delta}$ is a Markov transition probability matrix with matrix norm $\|E_{\delta}\|=1$, which implies that equation (\ref{eq:policy_valuation})
has a unique solution given by $V_{\delta_k}=[I-\beta E_{\delta_k}]^{-1} r_{\delta_k}$. Given the value $V_{\delta_k}$ the  policy improvement step
calculates a new decision rule $\delta_{k+1}$ given by the right hand side of equation (\ref{eq:bellman_policy}) except with 
$V_{\delta_k}$ substituted for $V$. 
If $\delta_{k+1}=\delta_k$ then policy iteration has converged and it is not hard to show that 
$V_{\delta_k}=V$, the unique solution to Bellman's equation (\ref{eq:bellman}).  

\subsection{Randomization, function approximation, and the Curse of Dimensionality}
\label{section:curse}

Inference in DDC and IRL requires repeated solutions of the MDP to find a reward function whose implied optimal behavior
best approximates observed behavior according to different metrics of goodness of fit, such as maximizing a likelihood 
function.  It is already challenging to solve a single MDP when $|S|$ or $|A|$ is large since each SA 
step requires $O(|A||S|^2)$ arithmetic operations to compute the expectation and maximization operators 
in equation (\ref{eq:bellman}) for all $s \in S$. Though PI requires only a few iterations to converge and thus
can be faster for smaller problems when $\beta$ is near 1, the downside is the $O(|S|^3)$ operations
required to solve the linear system (\ref{eq:policy_valuation}) in each policy valuation step.
In the absence of sparsity or other ``special structure'' alternative methods are required when $|S|$ or $|A|$ is large.

Bellman called this  the {\it curse of dimensionality\/} because it arises 
when  discretization and interpolation are used to approximate solutions to MDPs with multiple continuous state variables. 
Approximating a continuous MDP with $d$ state variables using a finite grid of $N$ points in each dimension 
results in an approximate finite MDP with $|S|=O(N^d)$ states which grow exponentially in $d$. The curse of dimensionality is unavoidable \citep{CT:1989}:  MDPs with $d_s$ continuous state and $d_c$ continuous choice variables
have a worst-case complexity (minimum number of arithmetic operations to compute an $\varepsilon$-approximation to $V$) 
that is bounded below by a function that grows at rate  $1/[(1-\beta)\varepsilon)]^{2d_s+d_c}$, a  lower bound that holds for {\it any possible 
algorithm.\/}

It may be possible to break the curse of dimensionality for subclasses of MDPs with special structure, such as  DDCs 
where the choice set $A$ is finite. When $|A|$ is small, most of the work to solve the MDP is spent on 1) numerical integration to
compute conditional expectations $EV(s,a)=\int V(s')p(s'|s,a)$, and 2) interpolation/approximation needed to find $V(s')$ for $s' \in S$.
While randomization techniques such as Monte Carlo integration have been proposed to mitigate the curse of dimensionality \citep{rust:1997}, their general applicability remains debated \citep{bray:2022}.

Another way to break the curse of dimensionality is to approximate $V$ by a function $V_\phi$  
that depends on a small number of unknown parameters $\phi$ for which $V_\phi$ approximately
solves the Bellman equation. For example $V_\phi$ could be the linear combination
$V_\phi(s)=\sum_{k=1}^K \phi_k \rho_k(s)$ where $\{\rho_k\}$ are ``basis functions'' and $\phi=(\phi_1,\ldots,\phi_K)$ are
coefficients, or $V_\phi$ could be a nonlinear function of $\phi$ such as a neural network.
Neural networks are universal approximators, which might help break the curse of dimensionality of approximating $V$, 
at least for certain subclasses of sufficiently smooth functions as first demonstrated by \citet{Barron1993}.
Let $V_{\hat\phi}$ denote 
a best approximation to $V$, where $\hat\phi$  
minimizes the mean squared ``Bellman residuals'' $\mbox{BE}(\phi)$
 over some grid of $N$ points $S_N=(s_1,\ldots,s_N)$ in $S$
\begin{equation}
      \hat\phi = \argmin_{\phi} \mbox{BE}(\phi)=E_{s \sim S_N}\left\{\left[V_\phi(s)- \Gamma(V_\phi)(s)\right]^2\right\},
\label{eq:nlls}
\end{equation}
where $E_{s \sim S_N}\{h(s)\}=\frac{1}{N} \sum_{i=1}^N h(s_i)$ denotes the sample average over the $N$ points in $S_N$.
Neural net approximation breaks the curse of dimensionality by allowing us to approximate $V$ with a small number
of network weights/parameters $\phi$.  Randomization also helps break the curse of dimensionality via
Monte Carlo integration instead of deterministic quadrature to approximate the  
conditional expectations $EV(s,a)$ entering the Bellman operator
$\Gamma(V_\phi)(s)$, transforming the solution of the Bellman equation for $V$ into a nonlinear regression problem (\ref{eq:nlls}).

Using the Triangle inequality and the contraction property of $\Gamma$, we can derive the error bound
$\|V-V_\phi\|_\infty \le \| V_\phi-\Gamma(V_\phi)\|_\infty/(1-\beta)$.\footnote{\footnotesize Here $\|f\|_\infty$
denotes the supremum norm of the function $f$, $\|f\|_\infty=\sup_{s \in S}\|f(s)\|$.} Thus, using a richer class of functions 
$\mathcal{F}_K=\{V_\phi|\phi \in R^K\}$ (i.e. a larger $K$) we can make the Bellman error
$\|V_\phi-\Gamma(V_\phi)\|_\infty$ smaller and better approximate the true solution to the Bellman
equation $V=\Gamma(V)$. In general, we require both $N$ and $K$ to increase to infinity at appropriate rates
to manage a classic bias-variance tradeoff to ensure that the nonlinear least squares estimator $V_{\hat\phi}$ (\ref{eq:nlls})
is a consistent estimator of the true $V$. Even when we can approximate $V$ well with relatively few parameters
$K$, there may still be a curse of dimensionality associated with finding a global minimizer $\hat\phi$.\footnote{\footnotesize 
Though stochastic optimization methods such as stochastic gradient descent can be helpful in avoiding getting stuck
at local minima, randomization does not break the curse of dimensionality of continuous non-convex optimization \citep{yn:1976}.}
Further, when $\beta$ is close to 1 $V_{\hat\phi}$ is not guaranteed to be close to the
true $V$. For example, if we start with 
the true solution $V=\Gamma(V)$ and perturb it by adding a constant $C$, the perturbed solution has a Bellman error of
$(1-\beta)C$ which can be small even when $C$ is large if $\beta$ is close to 1.
Fortunately, DDCs do not require
precise approximation of the level $V$ to accurately approximate the optimal policy in  (\ref{eq:bellman_policy}) since shifting $V$
by a constant leaves $\delta^*$ unaffected.\footnote{\footnotesize \citet{fujimoto22a} show that Bellman error is a poor proxy for value error in offline RL. 
\citet{Nguyen:2025} shows that in DDC problems, we obtain
more accurate solutions for $\delta^*$ by minimizing the sum of squared {\it anchored\/} Bellman errors, 
$V_\phi(s_j)-V_\phi(s_1)-[\Gamma(V_\phi)(s_j)-\Gamma(V_\phi)(s_1)]$ that difference out
parallel shifts in $V_\phi$, allowing the minimizer to more accurately approximate the {\it shape\/} of $V$ by ignoring its level, 
which is a version of ``relative value iteration'' proposed by \citet{puterman:1994} for solving MDPs under the long run average criterion.}


\subsection{Reinforcement Learning}
\label{section:rl}

These are stochastic iterative algorithms for approximating the value function $V$ and the optimal
policy $\pi$ inspired by models of trial and error learning by animals and 
humans.\footnote{\footnotesize  See \citet{BartoSutton:2020} and \citet{bertsekas:2025}
for excellent introductions to this literature.}
One of the first RL methods was 
{\it temporal difference learning\/} (TD) \citep{sutton:1988} which can be regarded as a type of asynchronous 
{\it stochastic approximation\/} \citep{robbinsmunro:1951} for solving the linear system (\ref{eq:policy_valuation}) for
$V_\delta$ in the policy valuation step of PI. TD is a special case of  
{\it Q-learning\/} \citep{Watkins:1989} where $Q$ is the {\it choice-specific value function\/} implied by the  Bellman equation
\begin{equation}
 \label{eq:q1}  
      Q(s,a)  =  r(s,a) + \beta \sum_{s'} V(s')p(s'|s,a)
  = r(s,a)+\beta \sum_{s'} \max_{a' \in A(s')} Q(s',a')p(s'|s,a) \equiv \Lambda(Q)(s,a).
\end{equation}
$Q$-learning can be viewed as a stochastic version of successive approximations 
$\{Q_t\}$ for finding the fixed point $Q=\Lambda(Q)$ via the updating rule
\begin{equation}
   Q_{t+1}(s_t,a_t)= Q_t(s_t,a_t) + \alpha_t \left[ r(s_t,a_t)+\beta \max_{a \in A(s_t)}[Q_t(\tilde s_{t+1},a)] - Q_t(s_t,a_t) \right],
\label{eq:qlearning}
\end{equation}
where $(s_t,a_t)$ is the state and action taken at  $t$, and $\tilde s_{t+1}$ is a random draw from the transition probability
$p(s'|s_t,a_t)$. \citet{Tsitsiklis:1994}
showed that if the step size sequence $\{\alpha_t\}$ converges to 0 at a sufficiently slow rate, the sequence $\{Q_t\}$ converges with probability 1 
to the unique $Q$ that solves (\ref{eq:q1}).\footnote{\footnotesize  \label{fn:step_size}
The conditions on the step size are 1) $\sum_t \alpha_t = \infty$
and 2) $\sum_t \alpha_t^2 < \infty$. Note $\alpha_t=1/t$ satisfies these conditions.}
Equation (\ref{eq:qlearning}) does not specify how  $a_t$ is chosen: optimality requires a {\it greedy policy\/}
$a_t = \argmax_{a \in A(s_t)} Q_t(s_t,a)$, but a key requirement for the convergence of Q-learning is 
that all pairs $(s,a)$ are sampled infinitely often. This presents an exploration/exploitation 
dilemma that implies that suboptimal actions must
be taken infinitely often in order for $Q$ learning to converge, but this
cannot occur under a greedy policy which is generically a pure strategy.  
One solution is to use an {\it $\varepsilon$-greedy policy\/} that takes the greedy action $a_t$
 with probability $1-\varepsilon$, and with probability $\varepsilon$ a randomly selected 
suboptimal action. However unless $\varepsilon$ converges
to $0$ with $t$, the policy we use to learn $Q$ is necessarily suboptimal and different from the optimal policy  we are trying to learn.
Thus Q-learning is an example of what \citet{BartoSutton:2020} refer  to as  {\it off-policy learning.\/}

\citet{softqlearning:2017} resolved the exploration/exploitation tradeoff  by 
 building on the {\it maximum entropy reinforcement learning\/} framework of \citep{ziebart:2008},
resulting in an algorithm they call {\it soft-Q learning.\/} 
They used Q-learning to solve a modified  MDP where the reward
 includes an entropy term that forces the optimal policy to be mixed rather than pure.
This ensures that Q-learning results in continuous exploration with all feasible actions being
chosen infinitely often, and they showed
it converges with probability 1 to the optimal mixed strategy $\pi$ given by the classic logit
or softmax formula that we will derive in equation (\ref{eq:ccp}) of section~\ref{sec:ddc}. 
The entropy of $\pi$ is $H(\pi(s))=-\sum_{a \in A(s)}\pi(a|s)\log(\pi(a|s))$. Adding it to the
reward results in the  MDP
\begin{equation}
       V(s)= \max_\pi  V_\pi(s),  \quad \mbox{where} \;  V_\pi(s) = E_\pi \left\{\sum_{t=0}^\infty \beta^t[r(s_t,a_t)+\sigma H(\pi(s_t))]| s_0=s\right\}.
\label{eq:vdefentropy}
\end{equation}
When $\sigma=0$ the objective (\ref{eq:vdefentropy}) reduces to the original MDP objective
(\ref{eq:vdef}) where $\pi$ is the pure strategy (\ref{eq:bellman_policy}). 
As $\sigma \to \infty$ the optimal policy for (\ref{eq:vdefentropy}) converges to a uniform distribution,
the maximum entropy solution. For $0 < \sigma < \infty$ the optimal
policy $\pi$ takes the form of a multinomial logit or ``softmax'' function (\ref{eq:ccp})
and the optimal value function $V$  in equation  
(\ref{eq:vdefentropy}) satisfies a ``soft Bellman equation'' identical
to the ``smoothed Bellman equation'' derived by \citet{rust:1988} for a ``partially observed'' MDP 
in equations  (\ref{eq:smoothV}) and (\ref{eq:smoothQ}) of section~\ref{sec:ddc}. In DDC models, individuals
use optimal pure strategies and the
entropy term $\sigma H(\pi(s))$ equals the expectation of components of rewards  that individuals
observe but we  do not observe. In contrast, the IRL literature uses entropy 
to enforce a direct preference for using mixed rather than 
pure strategies.\footnote{\footnotesize From an economic standpoint, it is not clear why individuals
should care about the entropy of a policy $\pi$: instead the random utility framework posits that
individuals care about the policy they use only to the extent that it directly maximizes their rewards.  
\citet{matejka:2015} provide an interpretation of the random utility model as a Bayesian decision problem
that involves a first stage choice of how much information to acquire about a finite set of alternative
choice options, where entropy captures the cost of information. They show that under an optimal
information acquisition strategy, the second stage choice probabilities take the multinomial or
``softmax'' form (\ref{eq:ccp}) derived in section~\ref{sec:ddc}. However this is a one shot model
of information acquisition: extending it to  DDC models involving sequential 
choices requires a Bayesian MDP formulation that includes
a component of the state variable that captures the posterior distribution of the individual's beliefs about the rewards
from various choices which is gradually updated over time, reflecting ``learning by doing.'' A classic example is the
``multi-armed bandit problem'' \citet{gittins:1974}.}



We say that RL is \textit{model-free} if it does not require explicit specification or knowledge of $p(s'|s,a)$.
Otherwise it is \textit{model-based}.  Q-learning can be  {\it model-free,\/}  unlike SA or PI, to the
extent that it does not require a specification for the
transition probability $p(s'|s,a)$ and avoids numerical integration to update $V$ and $Q$. This is possible
when ``the environment'' generates new states  rather than a computer simulation based on an explicit 
mathematical model for $p(s'|s,a)$. In  Watkins'  example
of an animal learning to adapt and thrive, we do not need to know $p$ since the environment ``draws'' the successor states
$s_{t+1}$ needed to implement the update in $Q_t$ in equation (\ref{eq:qlearning}). When trained in this manner Q-learning is indeed model-free
and so are other RL algorithms that interact directly with the environment (or use historical transitions for training/estimation as we discuss
in section~\ref{sec:irl}).

In many situations Q-learning and related RL methods produce policies that
are too noisy and converge too slowly to be useful in 
real applications, especially when the number of states is large.  As \citet{JiangXie:2024} note, their
``decisions can lead to undesirable outcomes, especially at the early stages of learning when the algorithm has little knowledge of the environment. This is not a problem when the environment is a simulator, but can lead to serious consequences when people are part of the environment.''
As a result RL training is increasingly done in simulation environments with
sufficient computer power to rapidly conduct many thousands or millions of training iterations to produce a much better solution than could be achieved
from limited real-world experience.  The ability to simulate from the true $p(s'|s,a)$  (or a good approximation to it)
for a large number of iterations is  key to producing good solutions using RL. For this reason most of the successful recent applications
of RL are  model-based in the sense that they rely on a model $p(s'|s,a)$ to simulate outcomes in offline
training even while using model-free RL algorithms to train the $Q$ values.\footnote{\footnotesize The use of Q-learning
to train chess strategies is an interesting intermediate case. Although the functional form of $p(s'|s,a)$
for chess is unknown since this transition probability depends on the action of the opponent, Q-learning can be classified as 
model-free since it relies on the environment to produce the next state $s_{t+1}$ directly
via self-play, without explicitly modeling or using the game's rules to calculate $p(s'|s,a)$.
In contrast, an algorithm like AlphaZero that uses Monte Carlo tree search explicitly uses the
game rules for planning and can be classified as model-based.}  

Tabular (i.e. finite state/action) Q-learning, (\ref{eq:qlearning}), is far too slow for MDPs with large state or action spaces. 
A scalable solution to this problem is to use function approximation methods 
such as Deep Q-Networks (DQN), which are  deep neural networks with weights $\phi$ 
that parameterize the Q-function as $Q_\phi$ similar to the way $V$ is approximated by $V_\phi$ in (\ref{eq:nlls}).
Instead of iteratively updating Q-values, as in (\ref{eq:qlearning}), 
the goal is to find $\hat\phi$ that minimizes the mean squared Bellman error
\begin{equation} 
   \mbox{BE}(\phi) = E_{(s,a) \sim \mathcal{D}} \left\{\left[\Lambda(Q_\phi)(s,a)-Q_\phi(s,a)\right]^2\right\}
=E_{(s,a) \sim \mathcal{D}}\left\{\left[ r(s,a)+\beta EV_\phi(s,a) -Q_\phi(s,a)\right]^2\right\},
\label{eq:dqnideal}
\end{equation}
where $\mathcal{D}=\{(s_i,a_i)|i=1,\ldots,N\}$ is a deterministic or random grid of $(s,a)$ pairs and $EV_\phi(s,a)$ is the
``emax'' function $\sum_{s'} \max_{a' \in A(s')} Q_\phi(s',a')p(s'|s,a)$.
Minimizing $\mbox{BE}(\phi)$ can be a difficult, high dimensional
optimization problem and we cannot even calculate it if the transition probability
$p$ is unknown.  There is also a danger that overfitting $Q$ can produce solutions with
low Bellman error $\mbox{BE}(\hat\phi)$ but suboptimal implied policies $\pi_{\hat\phi}$ \citep{fu2019diagnosing}.
However there are ways to deal with these challenges and  there is universal agreement that $\mbox{BE}(\phi)$ is the correct
objective to minimize when using DQN and other parametric methods to find approximate solutions to large scale MDPs.

An important goal in RL  is to try to avoid numerical integration with respect to  $p(s'|s,a)$ 
entering $\mbox{BE}(\phi)$ to obtain model-free learning similar to tabular Q-learning, (\ref{eq:qlearning}). This
suggests modifying the $\mbox{BE}(\phi)$ criterion to choose to $\hat\phi$ to minimize mean squared {\it temporal difference error\/}
\begin{equation}
      \mbox{TD}(\phi) = \mathbb{E}_{(s,a,s')\sim\mathcal{D}} \left\{ \left[r(s,a) + \beta \max_{a' \in A(s')} Q_{\phi}(s', a') - Q_\phi(s, a) \right]^2 \right\}.
\label{eq:td_loss}
\end{equation}
Unfortunately, it is well known that the $\hat\phi$ that minimizes $\mbox{TD}(\phi)$ 
differs from the value that minimizes $\mbox{BE}(\phi)$, since it is easy to see that
\begin{equation}
         \mbox{TD}(\phi) = \mbox{BE}(\phi) + \beta^2 \mathbb{E}_{(s,a,s')\sim\mathcal{D}} \left\{ \left[ \max_{a' \in A(s')} Q_\phi(s',a')
- EV_{\phi}(s,a)\right]^2\right\},
\label{eq:BETDdiff}
\end{equation}
a classic ``bias-variance'' decomposition,  where the second term on the right hand side of (\ref{eq:BETDdiff}), $\mbox{VQ}(\phi)$, is the
variance of the TD error.  

We now describe three main  modifications to the TD criterion (\ref{eq:BETDdiff}) that are used in current applications
of RL to achieve the benefits of model-free
simulation and training while still consistently estimating  the true $Q$ and $V$ in  (\ref{eq:q1}).
The first,  proposed by \citet{munos:2008}, is to parametrize the emax function $EV$ using separate parameters $\psi$
as  $EV_\psi(s,a)$ and use (\ref{eq:BETDdiff}) to write the Bellman error as $\mbox{BE}(\phi,\psi)=\mbox{TD}(\phi)-\beta^2 \mbox{VQ}(\phi,\psi)$,
and solve the minimax problem 
$(\hat\phi,\hat\psi)=\argmin_\phi\argmax_\psi \mbox{BE}(\phi,\psi)$. They showed that 
this modification results in a consistent estimator of $Q$ and $V$, but with the downside that the resulting minimax problem may result
``in considerable additional computational burden''. 
A second approach is to modify $\mbox{TD}$ to include
a moving  ``target function'' for $Q$, resulting in an iterative approach to minimizing BE called
{\it fitted value iteration\/} (FVI) \citet{antos:2008}.  Suppose we have some estimate of $\phi$  at
iteration $t-1$, $\hat\phi_{t-1}$. FVI results in an updated estimate $\hat\phi_t=\argmin_\phi\mbox{TD}(\phi,\hat\phi_{t-1})$
where the modified TD criterion using the target function $Q_{\hat\phi_{t-1}}$ is 
\begin{equation}
 \mbox{TD}(\phi,\hat\phi_{t-1}) = \mathbb{E}_{(s,a,s')\sim\mathcal{D}} \left\{ \left[r(s,a) +
 \beta \max_{a' \in A(s')} Q_{\hat\phi_{t-1}}(s', a') - Q_\phi(s, a) \right]^2 \right\}.
\label{eq:dqn_target}
\end{equation}
The target function ``trick'' converts the original TD objective (\ref{eq:td_loss}) into a standard linear or nonlinear regression, 
and it is easy to see that by doing this, the variance term in (\ref{eq:BETDdiff}) no longer depends on $\phi$, so when 
$Q_{\hat\phi_{t-1}}$ is close to the true $Q$, $\mbox{BE}(\phi)$ and $\mbox{TD}(\phi,\hat\phi_{t-1})$ are both minimized
by the same value $\hat\phi_t$, and the implied $Q_{\hat\phi_t}$ will be close to the true $Q$. 
The third approach is referred to as a {\it semi-gradient method\/} (see section 9.3 of
\citet{BartoSutton:2020}) since it equals the $\hat\phi$ that solves 
(or approximately solves) the $K$ ``moment conditions''
\begin{equation}
           0 = E_{(s,a,s') \sim \mathcal{D}} \left\{ \left[r(s,a)+\beta  \max_{a' \in A(s')} Q_{\phi}(s', a')
 -Q_\phi(s,a)\right]\nabla_\phi Q_\phi(s,a)
\right\},
\label{eq:semigradient_objective}
\end{equation}
where $\nabla_\phi Q_\phi(s,a)$ is the gradient of $Q_\phi$ with respect to $\phi$.\footnote{\footnotesize
It is referred to as a ``semigradient'' since the expression on the right hand side
of (\ref{eq:semigradient_objective}) is not a full gradient of  $\mbox{TD}(\phi)$
with respect to $\phi$ in (\ref{eq:td_loss}), but only with respect to the last term, $Q_\phi(s,a)$.}
If $Q=Q_{\phi^*}$, then it is easy
to see that for each $(s,a)$ the conditional expectation of $r(s,a)+\beta  \max_{a' \in A(s')} Q_{\phi}(s', a')-Q_\phi(s,a)=0$,
and hence the moment conditions defining the semi-gradient estimator  hold at $\phi=\phi^*$.
In the case of a linear approximator, if $Q(s,a)=\vec{\rho}(s,a)^{\top}\phi^*$ where $\vec{\rho}=(\rho_1,\ldots,\rho_K)$, then 
the $\phi^*$ that solves (\ref{eq:semigradient_objective}) is given by the {\it least squares TD\/} formula in section 9.8 of \citet{BartoSutton:2020},
\begin{equation}
       \phi^*= \left[E_{(s,a,s') \sim \mathcal{D}}\left\{ \vec{\rho}(s,a)^{\top}[\vec{\rho}(s,a)-\beta \max_{a' \in A(s')}\vec{\rho}(s',a')]\right\}\right]^{-1}
 E_{(s,a) \sim \mathcal{D}}\left\{\vec{\rho}(s,a)^{\top} r(s,a)\right\}.
\label{eq:lstd}
\end{equation}
Training of parameters in RL and ML models is customarily done sequentially and incrementally using 
{\it stochastic gradient descent\/} (SGD) where the parameters $\phi$ are updated in lockstep as each
new observation $(s_{t+1},s_t,a_t)$ arrives in  a simulated trajectory rather than only periodically after
the estimation objective and its gradient is evaluated for the full sample (or batch) of simulated outcomes $\mathcal{D}$. 
SGD can be viewed as the application of stochastic approximation to find a parameter $\phi$ that sets the
expected value of a system of equations to zero. In an optimization context such as RL, this system might be
the expected gradient of the objective function we are trying to minimize. 
For example consider the linear semigradient estimator for $\phi$ which solves (\ref{eq:semigradient_objective}).
Instead of directly solving it with, say, Newton's method, SGD solves it via the iterations
\begin{equation}
           \phi_{t+1}=\phi_t + \alpha_t\left[r(s_t,a_t)+\beta \max_{a' \in A(s_{t+1})} Q_{\phi_t}(s_{t+1},a')-Q_{\phi_t}(s_t,a_t)\right]
\nabla_\phi Q_{\phi_t}(s_t,a_t),
\label{eq:sgd_linearsemigradient}
\end{equation}
where $\alpha_t > 0$ is a step size parameter that satisfies the same conditions
as for tabular Q-learning in footnote~\ref{fn:step_size}. \citet{melo:2008} showed that if we approximate $Q$
as a linear combination of $K$ features $\vec{\rho}(s,a)$, $Q_\phi(s,a)=\vec{\rho}(s,a)^{\top}\phi$, then iterations of the form 
(\ref{eq:sgd_linearsemigradient}) converge with probability 1  to a parameter $\phi^* \in R^K$ satisfying
$Q_{\phi^*}=P_K(\Lambda(Q_{\phi^*}))$ where $P_K$ is the orthogonal projection operator onto the linear subspace 
spanned by the $K$ functions $\vec{\rho}(s,a)$. If $Q$ is in this subspace, then $Q=Q_{\phi^*}$, 
$P_K\Lambda=\Lambda$, and  $\phi^*$ is given by (\ref{eq:lstd}).

Counterexamples due to \citet{Baird:1995} and \citet{TsitsiklisVanRoy:1997} show that linear semigradient TD learning methods
such as (\ref{eq:sgd_linearsemigradient}) can diverge unless special care is taken. \citet{BartoSutton:2020} refer to the problem
as the  {\it deadly triad\/}  that can arise when off-policy training is combined with function approximation
and a  ``bootstrapping target'' (i.e. the 
term  $\max_{a' \in A(s')} Q_{\phi_t}(s_{t+1},a')$ in (\ref{eq:sgd_linearsemigradient})) that
changes too rapidly over iterations $t$. \citet{melo:2008} prove convergence under  
a condition on the moment matrix $E_\pi\{\vec{\rho}(s,a)^{\top}\vec{\rho}(s,a)\}$ of the basis functions
$\vec{\rho}(s,a)$ that guarantees that ``the trajectories of the algorithm to closely follow those of an associated
ODE with a globally asymptotically stable equilibrium point.'' (p. 666). 

However the main approach to ensuring stability of iterations with nonlinear approximators such as DQN is to use a ``target network''
that changes more slowly across iterations $t$.  \citet{antos:2008} provided sufficient conditions for convergence in probability 
of FVI iterations $\hat\phi_t =\argmin_\phi \mbox{TD}(\phi,\hat\phi_{t-1})$ where $\hat\phi_{t-1}$
constitute the target network parameters. However their analysis assumes  $\hat\phi_t$ is computed
in ``batch mode'' (i.e. using the entire simulated dataset $\mathcal{D}$) rather than updated recursively using
SGD.  \citet{zhang:2019} showed that SGD is convergent for linear approximators under ``a two-timescale framework, where the
main network is updated faster than the target network.'' \citet{Zhang:2023} provided sufficient conditions for
convergence of a SGD implementation of DQN, i.e. FVI using deep neural network  approximators $Q_\phi$ but using 
SGD iterations for $\phi_t$. Their analysis also 
allows for {\it experience replay,\/} i.e. the innovation of storing simulated outcomes in a replay buffer to help improve
stability and reduce the computational burden of SGD by using random sampling of simulated outcomes in ``minibatches'' 
from the replay buffer to reduce instabilities due to serially correlated training data and  smooth the learning process.

The DQN algorithm of \citet{Zhang:2023} involves a double loop of iterations, with $t=0,\ldots,T$ indexing the ``outer''
FVI iterations and $m=1,\ldots,M$ indexing the ``inner'' SGD iterations. Let $\phi_{t,m}$ be the parameter
value at iteration $(t,m)$ of this algorithm. The target network parameters are set to $\phi_{t,0}$ and are only
updated in the outer $t$ loop but are held fixed for all iterations $m$ in the inner SGD loop, whereas the inner
SGD iterations at FVI iteration $t$ are given by
\begin{equation}
     \phi_{t,m+1}=\phi_{t,m}+\alpha_t \left[\sum_{n \in \mathcal{D}_{t,m}} \left[r(s_n,a_n)+\beta \max_{a' \in A(s_{n+1})}
Q_{\phi_{t,0}}(s_{n+1},a')-Q_{\phi_{t,m}}(s_n,a_n)\right]\nabla_\phi Q_{\phi_{t,m}}(s_n,a_n)\right],
\label{eq:dqn_algorithm}
\end{equation}
where $\mathcal{D}_{t,m}$ is a {\it minibatch\/} of simulated trajectories drawn at random from 
the experience replay buffer $\mathcal{D}_t$ of trajectories $\{(s_n,a_n)|n=1,\ldots,N\}$ simulated for $N$ time steps
under a $\varepsilon_t$-greedy policy. The experience replay buffer $\mathcal{D}_t$ is refreshed at each step $t$ of the outer FVI
loop of the training process, but held fixed over the $M$ inner SGD iterations. Theorem 1 of \citet{Zhang:2023}
provides sufficient conditions for the DQN algorithm (\ref{eq:dqn_algorithm}) to generate estimates of $Q$ that
converge uniformly at rate $1/\sqrt{N}$, the size of the experience replay buffer: $\|Q_{\phi_{t,0}}-Q\|_\infty=O_p(1/\sqrt{N})$.

As a result of these algorithm design features that circumvent the deadly triad, DQN is not only convergent 
but has been successfully applied to solve a number of important large scale problems that were previously intractable.
\citet{mnih:2015} used a DQN to train effective strategies
for classic Atari 2600 video games. Similar methods were subsequently used for
board games such as chess, Go and Shogu, see \citet{Silver:2018}. In these applications, the transition probability
$p$ for the states is unknown due to the unknown response of opposing players. However 
they obtained a model-free implementation of DQN  by simulating game paths using {\it self-play\/} where
opponent players make choices according to the same $\varepsilon_t$-greedy policy used by the player
who is being trained. \citet{Zhang:2023} cites additional empirical successes of DQN in robotics and autonomous
vehicles.


%\subsection{Policy Gradient Methods}
%
%\textit{policy gradient} methods bypass solving the Bellman equation and directly learn a parameterized policy $\pi_\phi(a|s)$ directly by performing gradient ascent on the expected reward objective. The cornerstone of this approach is the Policy Gradient Theorem, which connects the performance gradient to the policy's score function, $\nabla_\phi \log \pi_\phi(a_t|s_t)$, weighted by a measure of subsequent reward. The earliest such algorithm, REINFORCE \citep{williams:1992}, uses the total Monte Carlo return of an episode for this weight. While this provides an unbiased gradient estimate, it suffers from high variance, as every action in a trajectory is credited with the same, often noisy, total reward. The standard solution is to subtract a state-dependent baseline, $b(s_t)$, from the return to reduce variance. The optimal baseline is the state-value function, $V_{\pi_\phi}(s_t)$, which leads to using the 
%\textit{Advantage Function}, $A_{\pi_\phi}(s_t, a_t) = Q_{\pi_\phi}(s_t, a_t) - V_{\pi_\phi}(s_t)$, as the learning signal \citep{kakade:2002}. This motivates the \textit{Actor-Critic} architecture, where a `critic' (a learned value function) provides a low-variance advantage estimate to guide the updates of the `actor' (the policy).
%
%While Actor-Critic methods address the variance problem, they introduce a stability challenge: a single large policy update, driven by an inaccurate advantage estimate, can lead to a catastrophic drop in performance. Subsequent research focused on ensuring stable, monotonic policy improvement. \textit{Trust Region Policy Optimization (TRPO)} \citep{schulman:2015} provided a theoretical solution by constraining the KL divergence between policy updates, but its implementation is computationally complex. Its successor, \textit{Proximal Policy Optimization (PPO)} \citep{schulman:2017}, achieves similar stability with a much simpler first-order algorithm using a novel clipped surrogate objective that penalizes large changes in the policy's probability ratios. PPO's balance of stability and implementation simplicity has made it a default algorithm, though its empirical success is now understood to depend not just on the clipping objective, but also on a suite of code-level optimizations like reward normalization and value function clipping \citep{engstrom:2019}. An alternative approach to stability is \textit{Soft Actor-Critic (SAC)} \citep{sac:2018}, which combines deep Q learning and policy gradient methods to
%minimize the Kullback-Leibler distance between $\pi_\phi$ and the softmax policy (\ref{eq:ccp}) implied by $Q_{\pi_\phi}$,
% encouraging exploration while stabilizing learning. Unlike PPO, SAC is off-policy and uses experience replay, making it more sample-efficient in continuous control tasks. PPO is a key component in training large language models via Reinforcement Learning from Human Feedback (RLHF), as we discuss in Section~\ref{sec:rlhf}.
%
%
 
